\section*{附录\quad 相关研究进展}
\addcontentsline{toc}{section}{附录\quad 相关研究进展}

本课题围绕动态负载下计算图的逐边执行机制选择展开：对计算图上的每一条连边，决定相邻两个计算阶段是在同一 persistent grid 内衔接、拆分为两次独立 launch 并按完成语义（completion）串行推进，还是在拆分边界上借助 Programmatic Dependent Launch（PDL）使后继提前进入并与前驱重叠；同时联合决定与之耦合的分组级资源配置（线程块数、分块形状、共享内存与寄存器预算）。当输入规模动态变化时，上述边界机制与资源配置的最优组合亦随之改变。围绕这一核心问题，近年文献可归纳为十三条研究线索：驻留式执行的演进与早期反思、大模型推理中的全图驻留执行、执行边界与后端的联合决策、融合编译器中的边界判据、动态 shape 与负载自适应优化、跨 kernel 重叠机制、kernel 内分工与跨 kernel 空间划分、GPU 片上资源共享的体系结构研究、分块与波量化（wave quantization）、性能代价模型、通信—计算重叠边、智能体驱动的 kernel 自动生成，以及推理服务系统与高性能算子库基线。其中前八条线索主要刻画执行形态的演进，后续线索则聚焦给定执行形态下的参数寻优与判据建模。

\subsubsection*{一、驻留式执行的演进与全图融合的固有代价}

将多个计算阶段放入单个长期驻留的 kernel 并非全新范式。\rcite{gupta2012}{Gupta 等 (2012)} 较早系统总结了 persistent threads 编程模型，通过启动恰好填满设备驻留容量的线程块并在设备端循环拉取任务，规避主机侧逐次 launch 开销及硬件调度器的不可控性。\rcite{steinberger2014}{Steinberger 等 (2014)} 提出的 Whippletree 在此基础上构建了面向动态异构任务的设备端调度框架。进入深度学习领域后，\rcite{diamos2016}{Diamos 等 (2016)} 的 Persistent RNN 将循环层权重常驻于片上存储以消除时间步间的重复访存，\rcite{holmes2019}{Holmes 等 (2019)} 的 GRNN 进一步将该思路拓展至低时延 RNN 推理。细粒度设备端运行时亦沿此路线演进：\rcite{yeh2017}{Yeh 等 (2017)} 的 Pagoda 利用常驻 MasterKernel 虚拟化 GPU 资源，按 warp 粒度调度并行度不足的窄任务，相对 CUDA Hyper-Q 取得平均 1.51$\times$ 加速；\rcite{belviranli2018}{Belviranli 等 (2018)} 的 Juggler 则在设备端完成依赖解析与任务派发，避免粗粒度全局屏障同步，在科学计算负载上相对全局屏障方案最高提升 31\%。

然而，“将完整计算图无条件融进单一 megakernel”在图形学领域很早便被证实存在固有代价。\rcite{laine2013}{Laine 等 (2013)} 在波前路径追踪研究中明确指出，单一 megakernel 的寄存器分配上限受其中资源需求最苛刻的阶段支配，其余轻量阶段被迫在同一低占用率（occupancy）约束下运行；将其拆分为若干独立阶段 kernel 并在阶段间重排任务反而显著更快。这一机理与本课题在 H100 上观测到的现象高度一致：单一 persistent kernel 的共享内存 carve-out 档位与每线程寄存器预算完全由最重阶段（如占用 255 个寄存器的 FlashAttention-2）决定，从而压低其余矩阵阶段的驻留上限。本课题在早期构建上的对照实验表明，当将 decode 阶段的 attention 从 255 寄存器实现替换为 91 寄存器实现后，物理切分相对最优单一 grid 的收益便由 14.4\%--16.5\% 收缩至 $-2.3\%$ 到 $+8.0\%$，印证了相当一部分切分收益源自解除重阶段对共驻代码的资源拖累。与路径追踪中动态数据驱动队列不同，大语言模型推理各阶段之间具有确定的数据依赖图结构，因而能够以算子边为单位精细权衡共驻与切分。

在深度学习编译器方向，\rcite{ma2020}{Ma 等 (2020)} 的 Rammer 提出 rTask 与 vEU 抽象，将算子内与算子间并行纳入同一静态调度空间并在 persistent kernel 中执行；\rcite{zhuang2024}{Zhuang 等 (2024)} 的 MonoNN 进一步将完整推理网络编译为单体 kernel，以发掘逐算子优化无法触及的跨阶段复用空间。这两项工作均将单一 kernel 视为既定终态，调度自由度仅存在于 kernel 内部，物理边界本身并非决策变量。与之互补，\cGraph 通过将多 kernel 序列固化为可重放有向无环图来摊销主机侧提交代价，但默认依赖边仍遵循“上游完全结束—下游方才启动”的串行完成语义。

\subsubsection*{二、大模型推理中的全图驻留执行}

近两年来，全图驻留式 megakernel 在低时延 LLM 推理领域得到集中推进。\rcite{spector2025}{Spector 等 (2025b)}（Hazy Research）面向 Llama-1B 手工构建了设备端任务图解释器，利用细粒度原子计数器替代全 grid 屏障维护跨阶段依赖，成为公开单请求 decode 时延最低的实现之一（本课题在 H100 上实测其 $B{=}1, S{=}512$ 单步时延为 1.08\,ms）。\rcite{he2026}{He 等 (2026)} 的 VDCores 针对异步 GPU 架构提出资源与执行解耦的编程模型；\cMPK 将设备端任务图方法系统编译器化，自动将张量程序降低为 SM 级任务与事件图并扩展至多卡通信场景，报告相对传统服务系统最高取得 1.7$\times$ 的端到端时延收益；\rcite{jin2026}{Jin 等 (2026)} 的 Event Tensor 则将分块任务间的依赖关系提升为一等张量抽象，通过编译器 ETC 联合实施静态与动态调度变换以支持动态 shape，并重点优化了固定形态 megakernel 在在线服务中的编译预热开销。

围绕上述主干路线近期出现了三方面扩展。其一是受限片上资源下的精细管理：\rcite{dong2026}{Dong 等 (2026)} 的 Ada-MK 针对消费级/推理级架构严格的共享内存与寄存器限制，借助 MLIR 细粒度依赖分解与离线静态轨迹搜索固化执行计划，并通过共享内存分页复用在在线广告推理场景中稳定超越 TensorRT-LLM。其二是基于智能体的自动代码合成：\rcite{li2026f}{Li 等 (2026a)} 的 ForgeMegakernel 结合里程碑知识库与中间态校验预言机自动生成 decode megakernel，在 0.6B--13B 模型上报告了相对 SGLang 平均 1.21$\times$、相对 MPK 平均 1.54$\times$ 的加速；\rcite{jaber2026}{Jaber 和 Jaber (2026)} 的 AutoMegaKernel 引入静态调度 IR 验证器以在启动前排除死锁与数据竞争，但也坦陈在 A100/H100 上受跨 SM 同步开销制约，部分场景仍慢于经 CUDA Graph 封装的 cuBLAS。其三是面向异构计算—通信混合负载的核内空间划分：\rcite{huang2026}{Huang 等 (2026a)} 的 Weave 与 \rcite{mao2026}{Mao 等 (2026)} 的 mKernel 在单核或跨机融合 kernel 内部将 SM 动态划分为计算组与通信组，分别在多卡 MoE 与 Ring Attention 等负载上取得显著重叠收益。此外，\rcite{yangy2026}{Yang 等 (2026)} 的 GPUOS 探索向运行中的常驻 kernel 动态注入即时编译算子，\rcite{chowdhary2026}{Chowdhary 等 (2026)} 的 Fleet 引入面向多芯粒（multi-die）GPU 的局部 L2 亲和任务层级，而 \rcite{guo2026}{Guo 和 Miao (2026)} 的 PipeLink 则利用资源契约静态检验加载、主循环与收尾组件能否安全链接进同一 kernel。

总体来看，这一方向仍普遍以单一驻留 kernel 作为默认目标形态，将优化重心放在 kernel 内部的 SM 角色划分与事件同步上，因而未能从根本上摆脱异构阶段共驻时的包络木桶效应。

\subsubsection*{三、执行边界与后端的联合决策}

近期若干前沿工作开始正面审视“在哪里切分 megakernel”这一问题。\cChat 将流式推理模型的分段切分显式建模为连续区段划分、区段后端选择（常规算子图或 megakernel）及跨区段共享策略的三元组搜索，依托独立算子剖析与排序代理评分免去端到端遍历，发现 NH-TTS、Qwen3-TTS 与 Llama-1B 上的最优方案随 batch 规模在零至四个 megakernel 之间动态迁移；\rcite{chatterjee2026b}{Chatterjee 等 (2026b)} 进一步在连续批处理迭代边界上根据实时并发度在整图 megakernel 与常规融合链之间二选一，在 128 并发与 8 并发区间分别兼顾了高吞吐与低时延。

与此同时，\cQiang 从微观机制层给出了极具启发性的实证反思。该工作基于同一份底层算子代码将 ThunderKittens 与 VDCores 两个代表性 megakernel 拆解重构为独立 kernel 序列，证实线程块常驻并不天然等价于有效计算重叠：配合 PDL 重叠启动与权重加载后，拆分版 ThunderKittens 的 decode 时延反而降低 5.7\%。作者通过细粒度插桩指出，设备端细粒度发布与内存序维护本身存在可观开销，在途异步访存流量较高时信号发布时延显著上升；更重要的是，persistent 线程块可在内部直接复用片上存储，而分离 launch 的后继 kernel 必须在 SM 上获取完整的线程块资源配额后方能准入执行。不过，该文将这一准入限制归结为对未来 GPU 细粒度跨核资源共享硬件支持的诉求，尚未在现有硬件上提炼出可定量预测的准入判据。本课题前期工作 DACOS（\cDacos）则从固定算子边界出发，在合法性约束下联合决策后继预启动时机、操作数预取规模与存储层级放置。

此外，相邻方向还在空间流水与编译搜索层面展开了探索：\rcite{andoorveedu2026}{Andoorveedu 等 (2026)} 的 SubGPU Pipelining 与 \rcite{davies2025}{Davies 等 (2025)} 的 Kitsune 将计算图切分为绑定于互斥 SM 子集的空间流水段；\rcite{dai2025}{Dai 等 (2025)} 的 STOF 与 \rcite{huangz2026f}{Huang 等 (2026b)} 的 FlashFuser 联合搜索算子融合边界与分块流水配置；\rcite{lix2026}{Li 等 (2026c)} 的 AgentCompile 在四档离散融合粒度间择优；\rcite{linj2026}{Lin 等 (2026b)} 的 MegaOps 则在保持 kernel 分离的前提下将依赖粒度下沉至单个中间分块。在工业界实践中，NVIDIA Transformer Engine 的 Userbuffers 后端（\rcite{nvidiate}{NVIDIA, 2026g}）、CoCoNet（\rcite{jangda2022}{Jangda 等, 2022}）与 PyTorch async-TP（\rcite{wangy2024}{Wang 等, 2024}）已支持逐算子静态配置重叠策略；\rcite{together2026}{Together AI (2026)} 的工程报告更直接指出，H100 上 megakernel 内部全局同步（约 646\,ns）与 CUDA Graph 配合 PDL 的跨边界交接（约 645\,ns）开销几乎一致，迫使全图统一 warp 数与 cluster 规模反而带来显著减速，因此转而为模型各条边手工标注同步方式。

对比可见，现有工作或以粗粒度区段与整次前向为单位切换后端，或仅在离线初始化阶段静态指定重叠策略，尚未将“同 grid 共驻、跨 grid completion 与跨 grid PDL”作为逐边决策变量与分组级资源配置联合求解，亦缺乏在现役硬件上直接判定 PDL 准入与落位代价的解析依据。

\subsubsection*{四、融合编译器中的边界判据}

传统深度学习图编译器围绕算子融合边界积累了四类代表性判据。第一类是基于算子属性分类的启发式规则，如 TVM（\rcite{chen2018}{Chen 等, 2018}）、Ansor（\rcite{zhengl2020}{Zheng 等, 2020}）、DNNFusion（\rcite{niu2021}{Niu 等, 2021}）与 Hidet（\rcite{ding2023}{Ding 等, 2023}）。第二类是基于访存流量与并行度衰减的解析代价模型，如基于 tile-graph 的 Welder（\rcite{shi2023}{Shi 等, 2023}）、Chimera（\rcite{zhengs2023}{Zheng 等, 2023}）、MCFuser（\rcite{zhangzh2024}{Zhang 等, 2024}）及分层反馈融合框架 Apollo（\rcite{zhao2022}{Zhao 等, 2022}）。第三类是精确数学规划或超优化枚举，如将原语重组表述为 0-1 线性规划的 Korch（\rcite{hu2024}{Hu 等, 2024}）与多级超优化器 Mirage（\rcite{wu2025}{Wu 等, 2025}）。第四类则如 Triton（\rcite{tillet2019}{Tillet 等, 2019}）与 TileLang（\rcite{wangl2025}{Wang 等, 2025}），由程序员显式界定 kernel 范围而由编译器优化内部流水。

上述工作中有两项设计与本课题密切相关。一方面，片上驻留资源上限已被引入融合切分判据：例如 Souffle（\rcite{xia2024}{Xia 等, 2024}）要求子程序内部全局同步所需的最大占用不超过单波驻留容量，超限即切分为新 kernel；AStitch（\rcite{zhengz2022}{Zheng 等, 2022a}）则根据共享内存与寄存器压力在寄存器、共享内存和全局内存多级缝合方式间降级。但这些判据仅检验单个融合组自身的驻留合法性，并未刻画相邻两个独立 kernel 能否在同一 SM 上跨边界交叠驻留。另一方面，AStitch 与 Welder 虽支持逐边选择数据复用层级，但其候选操作均局限于同一 kernel 内部，并未纳入“保留物理边界并利用硬件 PDL 抵消交接开销”这一跨边界维度。

\subsubsection*{五、动态 shape 与负载自适应优化}

真实推理服务中请求到达率、提示长度与缓存上下文持续波动，使面向动态 shape 的编译优化成为本课题系统设计的重要参照。该领域的经典架构普遍遵循“离线构建候选空间、在线轻量匹配”的分工模式：Nimble（\rcite{shen2021}{Shen 等, 2021}）、DISC（\rcite{zhu2021}{Zhu 等, 2021}）与 BladeDISC（\rcite{zhengz2024}{Zheng 等, 2024a}）在编译期保留符号化维度关系；DietCode（\rcite{zhengb2022}{Zheng 等, 2022b}）、MikPoly（\rcite{yu2024}{Yu 等, 2024}）与 MoonPoly（\rcite{zhangy2026}{Zhang 等, 2026a}）离线生成固定尺寸微内核并在运行时依实际输入规模动态拼合；FlashDecoding++（\rcite{hong2024}{Hong 等, 2024}）、WaveTune（\rcite{zhang2026}{Zhang 等, 2026b}）与 DeepGEMM（\rcite{deepgemm2025}{DeepSeek, 2025a}）则根据矩阵形状动态切换算子内部的分块配置或数据流。

然而，上述动态 shape 编译器与算子库在运行时调整的对象均限于单个算子内部的分块参数与微内核组合，算子间始终固定为逐 kernel 串行启动；而前述 megakernel 自适应工作（\cChat；\rcite{chatterjee2026b}{Chatterjee 等, 2026b}；\rcite{wu2026}{Wu 等, 2026}；\rcite{yadav2026}{Yadav 和 Zhao, 2026}）又仅在整段或整次前向粒度上切换后端。本课题的跨负载实验（正文图~\ref{fig:motiv} 与第~\ref{sec:types}~节）表明，在无约束逐负载寻优下 15 个典型负载会产生 11 种不同的逐边方案；而进一步引入 1\% 时延容差作敏感度解耦后可以发现，边级切分机制在 decode 上可离线固定、在 prefill 上仅需两个协同候选即可覆盖，真正沿行数轴频繁更替的是各分组的分块形状与线程块配置（约需 22 个行数区间）。这一认识将动态 shape 自适应自然推广至“少量机制候选 + 分段资源配置查表”的联合空间。

\subsubsection*{六、跨 kernel 重叠执行机制}

作为 Hopper 架构引入的硬件级跨边界协同原语，PDL 允许前驱通过 \texttt{griddepcontrol.launch\_dependents} 提前释放后继调度资格，并由后继在首次依赖读取前通过 \texttt{cudaGridDependencySynchronize}（对应 PTX 指令 \texttt{griddepcontrol.wait}）保证数据可见性（\cPDL）。当前主流推理框架（如 SGLang 的低时延模式）多将 PDL 视为全局开关对前向链近乎全量开启；而 Hazy Megakernels（\rcite{spector2025}{Spector 等, 2025b}）则认为其同步粒度偏粗而完全退回单核内部同步。此外，SET（\rcite{li2026s}{Li 等, 2026b}）、Blink（\rcite{siavashi2026}{Siavashi 等, 2026a}）、设备端多卡任务图（\rcite{mustard2025}{Turimbetov 等, 2025}）以及体系结构层的 Wireframe（\rcite{wireframe2017}{Abdolrashidi 等, 2017}）与 BlockMaestro（\rcite{blockmaestro2021}{Abdolrashidi 等, 2021}）分别从流事件图重用、网卡/设备端自驱动与硬件依赖追踪角度压缩跨核调度间隙。相比将 PDL 当作经验性全局开关的既有做法，本课题首次从片上存储档位兼容性与线程块派发动力学角度解释了“为何同一机制在不同边与负载上有正有负”，并给出了可静态判定的选择准则。

\subsubsection*{七、kernel 内分工与跨 kernel 空间划分}

在单个 Hopper kernel 内部，以 CUTLASS ping-pong GEMM 为代表的实现广泛采用 warp 专用化（warp specialization），令生产者 warp 通过 TMA 异步填充共享内存、消费者 warp group 执行 WGMMA 并借助 \texttt{setmaxnreg} 非对称分配寄存器；公开优化指南虽充分探讨了流水级数与单核占用率的权衡，但较少涉及共享内存 carve-out 对跨 kernel 交接的连锁影响。在跨 kernel 空间复用方向，CUDA 官方提供的 MPS（\rcite{nvidiamps}{NVIDIA, 2026c}）、MIG（\rcite{nvidiamig}{NVIDIA, 2026d}）与 Green Context（\rcite{nvidiagreen}{NVIDIA, 2026e}）或缺乏硬性 SM 隔离与前向进展保证，或重配置粒度过粗且显存开销高昂；近期推理服务工作如 Bullet（\rcite{lin2026}{Lin 等, 2026a}）、Nova（\rcite{xu2025n}{Xu 等, 2025}）、DuetServe（\rcite{gao2025}{Gao 等, 2025}）与 SubGPU Pipelining（\rcite{andoorveedu2026}{Andoorveedu 等, 2026}）普遍因 SM 内并发性能不可控而退回 SM 集合间的空间切分，POD-Attention（\rcite{kamath2025}{Kamath 等, 2025}）及算子内融合工作（\rcite{jangda2026}{Jangda 等, 2026}）则重新将并发逻辑收归单一融合 kernel 内部。与上述面向无依赖任务划分 SM 归属的工作不同，本课题聚焦具有直接数据依赖的前后继分组在同一 SM 集合上的时序交叠规律。

\subsubsection*{八、GPU 片上资源共享与 PDL 准入的先行记录}

计算机体系结构领域长期关注 SM 内寄存器文件、共享内存与线程槽难以跨 kernel 细粒度共享的瓶颈，先后提出了共享内存多路复用（\rcite{yang2012}{Yang 等, 2012}；\rcite{jatala2017}{Jatala 等, 2017}）、寄存器虚拟化与分时复用（\rcite{jeon2015}{Jeon 等, 2015}；\rcite{khorasani2018a}{Khorasani 等, 2018a,b}；\rcite{liuz2026}{Liu 等, 2026b}）、统一资源虚拟化 Zorua（\rcite{vijaykumar2016}{Vijaykumar 等, 2016}）以及并发多核切片执行框架 Kernelet、SMK、Warped-Slicer、弹性 kernel 与 AMD 软件共驻模型（\rcite{adriaens2012}{Adriaens 等, 2012}；\rcite{pai2013}{Pai 等, 2013}；\rcite{kernelet2013}{Zhong 和 He, 2013}；\rcite{wangz2016}{Wang 等, 2016}；\rcite{xu2016}{Xu 等, 2016}；\rcite{wuh2020}{Wu 等, 2020}）。但上述方案绝大多数依赖硬件微架构改动，或仅面向互无数据依赖的多任务吞吐优化；NVIDIA 官方占用率指南（\rcite{nvidiabp}{NVIDIA, 2026h}）亦仅定义单一 kernel 的驻留上限。

梳理与本课题 PDL 准入判据直接相关的公开记录，可以看到四层渐进线索。首先，在官方编程指南层面，自 Volta 架构恢复 L1 与共享内存统一片上存储后，\rcite{nvidiacarve}{NVIDIA (2026i)} 指出切换共享内存容量配置需等待 SM 排空，并提示低共享内存 kernel 可主动请求较大 carve-out 以便与后续 kernel 并发，但未给出定量档位规则，PDL 文档更仅将提前启动表述为机会性行为。其次，在开发者社区答复中，NVIDIA 工程师针对多流并发落入互斥 SM 的现象解释称，Hopper 按单 SM 粒度动态调整 L1/共享内存配比，小档位 kernel 占据 SM 期间会阻塞大档位 kernel 共驻，反之则允许进入（\rcite{nvidiaforum}{NVIDIA, 2025}）。再次，Mosaic（\rcite{strati2026}{Strati 等, 2026}）与 FlashInfer 在 Blackwell 上的问题报告（\rcite{flashinfer7548}{FlashInfer, 2026}）在多流干扰场景中引述并复现了这一现象。最后，在 PDL 场景下，\cQiang 与 CUTLASS 社区讨论（\rcite{cutlass2778}{CUTLASS, 2025}）均只将准入失败归因于寄存器或共享内存余量不足，未触及 carve-out 档位约束。本课题首次将“carve-out 档位单调兼容与链式保持”同“SM 剩余资源容量”统一为可静态计算的 PDL 准入充要判据，并揭示了准入受阻且后继未满驻留时引发的集中放置开销与 L1 容量缩减效应。

\subsubsection*{九、分块、波量化与工作划分}

一旦允许各分组独立选择分块形状与线程块数，波量化（wave quantization）便成为决定阶段时延阶梯特征的核心因素：当输出分块总数不是设备单波驻留线程块数的整数倍时，尾波仅占用少量 SM 而耗时仍接近完整一波（\rcite{nvidiamm}{NVIDIA, 2026f}）。Hopper 架构拥有 132 个 SM（含素因子 11），使 2 的幂次分块更易遭遇非整波尾效应（\rcite{siavashi2026b}{Siavashi 等, 2026b}）。为此，Stream-K（\rcite{osama2023}{Osama 等, 2023}）与面向 decode attention 的 LeanAttention（\rcite{sanovar2024}{Sanovar 等, 2024}）采用以工作单元为中心的均匀切分消除尾波，TokenWeave（\rcite{gond2026}{Gond 等, 2026}）与 MoE 通信解析模型（\rcite{liuh2026}{Liu 等, 2026a}）则依据波量化结构反推切分比例或通信线程块数。本课题实测同样印证了这一规律：在 64 行输入、64 分块下，U 阶段采用 264 个线程块（恰好整波）耗时 28.8\,\textmu s，而改用 132 个或 396 个线程块时分别升至 34.0\,\textmu s 与 33.3\,\textmu s，相差达 18\%。这也从侧面说明，若单一 grid 基线被迫对全部阶段施加统一线程块数，其性能劣势在很大程度上源于波量化失配而非同步原语本身。

\subsubsection*{十、面向配置选择的性能代价模型}

为将全图规划从穷举实测推进至快速搜索，需要能够由离线算子特征外推分组与边界表现的性能模型。早期 GPU 解析模型如 Hong-Kim 模型（\rcite{hong2009}{Hong 和 Kim, 2009}）、GPUMech（\rcite{huang2014}{Huang 等, 2014}）、Accel-Sim（\rcite{khairy2020}{Khairy 等, 2020}）与 GCoM（\rcite{lee2022}{Lee 等, 2022}）均面向独占设备的单一 kernel。近期面向张量分块与大模型推理的工作显著提升了单核与单组预测精度：TileSight（\rcite{mo2026}{Mo 等, 2026}）在分块粒度建模访存复用与流水线向量，WaveTune（\rcite{zhang2026}{Zhang 等, 2026b}）构建波结构感知的双线性时延模型，ForeSight（\rcite{wang2026}{Wang 等, 2026}）通过还原硬件单元排队停顿将单核误差压至 8.3\%，Bullet（\rcite{lin2026}{Lin 等, 2026a}）与 KernelSight-LM（\rcite{yao2026}{Yao 等, 2026}）引入随 SM 份额伸缩的 roofline 模型，而 Welder、Chimera、MCFuser 与 Neptune（\rcite{zhao2025n}{Zhao 等, 2026}）则对融合子图整体估价。然而，现有模型均未显式刻画跨边界交接时的动态边代价——即 PDL 后继能否提前准入、准入后线程块落位偏斜及继承 L1 缩小带来的时延增减。为此，本课题结合设备端纳秒计时器与微基准探针，将准入、落位与预取三项效应逐一剥离量化，再与分组级波量化模型组合驱动规划。

\subsubsection*{十一、作为同类边决策的通信—计算重叠}

分布式训练与推理中的通信—计算重叠在抽象结构上与本课题高度同构，均可归结为相邻计算与通信阶段是否共驻融合、以何种粒度同步以及如何分配 SM 资源。现有方案大体沿三条技术路径展开：第一类采用粗粒度张量切分配合多流异步并发或拷贝引擎卸载，代表工作包括操作序列分解（\rcite{wangs2023}{Wang 等, 2023}）、MegaScale（\rcite{jiang2024}{Jiang 等, 2024}）、Centauri（\rcite{chen2024}{Chen 等, 2024}）、TokenWeave（\rcite{gond2026}{Gond 等, 2026}）及 PyTorch async-TP（\rcite{wangy2024}{Wang 等, 2024}）；第二类采用分块级深度融合，将集合通信原语直接编织进 GEMM 的前导或收尾阶段，如 CoCoNet（\rcite{jangda2022}{Jangda 等, 2022}）、FLUX（\rcite{chang2024}{Chang 等, 2024}）、TileLink（\rcite{zhengs2025b}{Zheng 等, 2025b}）、Triton-distributed（\rcite{zhengs2025a}{Zheng 等, 2025a}）、硬件跟踪机制 T3（\rcite{pati2024}{Pati 等, 2024}）、Syncopate（\rcite{qiang2026a}{Qiang 等, 2026a}）与 mKernel（\rcite{mao2026}{Mao 等, 2026}）；第三类则聚焦跨流并发时的显式 SM 配额约束与产出次序重排，如 Entwine（\rcite{mak2026}{Ma 等, 2026}）、Transformer Engine 与 DeepEP（\rcite{deepep2025}{DeepSeek, 2025b}）。已有实验普遍显示，通信融合或 SM 预留在小规模输入下易因寄存器膨胀或计算波数增加而出现负收益，且现有方案均采用离线静态配置并较少结合 PDL 机制，这为本课题后续将通信节点纳入动态逐边决策框架提供了直接切入点。

\subsubsection*{十二、智能体与搜索驱动的 kernel 自动生成}

近年来基于大语言模型智能体与强化学习的 CUDA 代码生成发展迅速，但以 KernelBench 及工业界大规模演进系统 KernelEvolve（\rcite{liao2025}{Liao 等, 2025}）为代表的主流工作仍以给定算子边界内的单核加速比为优化目标，跨算子边界划分与图级协同仍被列为开放问题；少数触及图级的代码生成框架（如 AgentCompile、ForgeMegakernel 与 AutoMegaKernel）也仅在粗粒度融合档位或整图单一 megakernel 上寻优，尚未覆盖逐边异构执行机制与动态资源包络的联合空间。

\subsubsection*{十三、推理服务系统与高性能算子库基线}

在系统层面，vLLM（\cVllm）、SGLang（\rcite{zhengl2024}{Zheng 等, 2024b}）与 TensorRT-LLM 代表了当前“逐算子高度优化 + 连续批处理 + CUDA Graph / 全局 PDL”的主流工业基线，NanoFlow（\rcite{zhuk2025}{Zhu 等, 2025}）进一步通过算子级 nano-batch 流水提升单卡吞吐。在算子实现层面，attention 沿着 FlashAttention-2（\cFA）、利用 Hopper 异步 TMA 与 warp 专用化的 FlashAttention-3（\rcite{shah2024}{Shah 等, 2024}）、面向解码并行归约的 Flash-Decoding（\rcite{dao2023}{Dao 等, 2023}）及可定制引擎 FlashInfer（\rcite{ye2025}{Ye 等, 2025}）持续迭代；矩阵与通用模板则有 ThunderKittens（\rcite{spector2025a}{Spector 等, 2025a}）与运行时即时编译的 DeepGEMM 等代表。这些系统共同构成了检验本课题方法有效性的外部参照，同时也要求本课题在统一底层算子代码的受控条件下剥离出纯属边处理机制与资源配置解耦的真实收益。

\subsubsection*{十四、小结：现有研究留下的关键空白}

综合以上十三条线索，现有研究在低时延动态推理执行计划上仍留下六点明确空白，构成本课题的核心切入点。第一，在决策粒度与候选机制上，既有工作或局限于整图、连续区段与单次前向切换，或仅在单 kernel 内部选择存储复用方式，尚未将“同 grid 共驻、跨 grid completion 与跨 grid PDL”作为逐边变量同分组级分块与线程块配置联合优化。第二，在动态输入规模自适应上，动态 shape 编译仅调节算子内部参数，而 megakernel 切换粒度过粗，缺乏对边级执行机制与资源配置随规模变化规律的解耦分析。第三，在跨边界重叠机理上，文献对 PDL 的正负收益多停留于经验观察或未来硬件诉求，缺乏结合共享内存 carve-out 档位与资源余量的现役硬件可预测判据。第四，在片上资源共享范式上，无需硬件改动的软件方法此前多面向无依赖并发任务或静态 SM 分区，尚未系统利用配置选择引导有依赖前后继 kernel 在同一 SM 上安全交叠。第五，在性能代价建模上，现有解析模型均只评估孤立 kernel 或融合子图，尚未纳入跨边界准入、落位偏斜与 L1 档位继承所构成的边级交互代价。第六，在公平评测方法学上，单核基线若缺乏分阶段驻留与逻辑块号重排等同等调优，极易夸大物理切分收益；本课题因此在完全一致的算子实现与充分调优的单一 grid 基线上开展归因研究。
